\subsection{Трассы}

Как правило, при разработке сложных микроархитектур вначале пишется не RTL-код, а потактовый симулятор разрабатываемого устройства на языке более высокого уровня(например, C++). Основное назначение такого симулятора --- составить референсную модель. С помощью этой рефенсной модели можно проанализировать происзводительность на сложных и длинных тестах, которые затрачивали бы нереаличтичные промежутки времени на потактовом симуляторе RTL. 

Вместе с тем, в процессе анализа формируются векторы входных с выходных данных, которые проходят через эту модель. Их будем называть трассами. Для проверки соответствия разработанного ртл модуля можем использовать эти трассы как средство верификации. 

Практически таким же образом можем использовать чекпоинты --- снимки состояния системы в определённые моменты времени. 

